\documentclass[12pt]{article}
\usepackage[utf8]{inputenc}
\usepackage[T1]{fontenc}
\usepackage{hyperref}
\title{Sened Language Resources for K-12 Curriculum Development}
\author{Compiled by Claude Sonnet 5.5 (claude-sonnet-5-5) for Andr\'as Kornai}
\date{}
\begin{document}
\maketitle

\section{Introduction}
Sened is an extinct East Zenati Berber (Afro-Asiatic) language that was spoken in the town of Sened and nearby Majoura (Tmagurt) in southern Tunisia. Sources disagree on when it died out. In 1911 the whole town still spoke Berber, but by 1968 only the elderly did. Wikipedia and ASJP give 1990 as the date of extinction. Kossmann places it in the Tunisian-Zuwara cluster, together with Matmata, Douiret, Jerba and Zuwara. This dossier rests on a single search. Further searches were not possible in this session because the search budget was exhausted. The documentation listed here is therefore very thin and not exhaustive.

\section{Primary Educational Resources}
No primers, textbooks or school programmes for Sened turned up in the search. A curriculum would have to be built from historical documentation (see below), and any teaching would be heritage-awareness and revival work rather than instruction in a living language.

\section{Linguistic and Reference Materials}
\begin{itemize}
\item Wikipedia article on the Sened language, with classification and extinction dates: \url{https://en.wikipedia.org/wiki/Sened_language}
\item Wikipedia article on the town of Sened, Tunisia: \url{https://en.wikipedia.org/wiki/Sened,_Tunisia}
\item ASJP database entry for Sened (a word list): \url{https://asjp.clld.org/languages/SENED}
\item Rosetta Panglossia (Long Now) entry for ISO code sds: \url{https://rosettapanglossia.longnow.org/index.php?title=Sds}
\item Wikipedia article on the East Zenati languages, which gives the wider classification: \url{https://Www.Wikipedia.org/wiki/East_Zenati}
\end{itemize}
The main documentation, according to the search summary, is Provotelle's 1911 \emph{\'Etude sur la Tamazir't, ou, Z\'enatia de Qala\^at Es-Sened}. It contains a short excerpt from the epic \emph{Taghribat Bani Hilal}. I did not find a URL for it in the search results.

\section{Religious and Community Resources}
No religious or community resources for Sened were found.

\section{Digital and Audio Resources}
No audio, video or software resources were found. The language was extinct before any recordings appear to have been made, so none are likely to exist.

\section{Curriculum Development Context}
Sened has no known speakers and almost no documentation, so a curriculum can use the language itself only to a very limited extent. The realistic approach is a heritage-awareness unit on the Berber history of southern Tunisia, built from the Provotelle text and the word list. Related living varieties (the Tunisian-Zuwara cluster) could supply comparative material. More searching, for instance for the Provotelle text and for Tunisian Berber revitalization groups, is recommended before relying on this dossier.

\end{document}